\usepackage[T1]{fontenc}
\usepackage{xcolor}
\usepackage{listings}
\usepackage{inconsolata}
\usepackage{relsize}
\usepackage{amsmath,amssymb,array}
\usepackage{esz}
\usepackage{mathpartir}
\usepackage{graphicx}
\usepackage{xspace}
\usepackage{xurl}
\usepackage{caption}
\usepackage{bbding}
\usepackage{tabularx}
\usepackage{enumitem}
\usepackage{cite}
\usepackage[compact]{titlesec}
% IEEEtran.cls sets \citepunct to "], [", which produces "[A], [B]" for
% \cite{A,B}. Override to a plain ", " so consecutive cites share one bracket.
\renewcommand{\citepunct}{, }

\hypersetup{colorlinks=true,linkcolor=blue,citecolor=blue,urlcolor=blue}
\def\UrlFont{\rmfamily}

% Tighten section/subsection vertical spacing for two-column layout.
\titlespacing*{\section}{0pt}{1.2ex plus 0.2ex minus 0.2ex}{0.6ex plus 0.1ex}
\titlespacing*{\subsection}{0pt}{0.8ex plus 0.1ex minus 0.1ex}{0.4ex plus 0.1ex}

\newcommand{\Section}[1]{\section{#1}}
\newcommand{\SubSection}[1]{\subsection{#1}}
\newcommand{\Paragraph}[1]{\vspace{0.25ex}\noindent\textbf{#1.}\xspace}

% Labeled-paragraph list: bold inline label, lightly indented.
% Use as \begin{Description} \item[Label.] body \item[...] ... \end{Description}.
\newlist{Description}{description}{1}
\setlist[Description]{leftmargin=1em, labelindent=1em, style=unboxed, font=\bfseries, itemsep=0.5ex, topsep=0.5ex}

% Bulleted list with an em-dash marker flush against the left text margin.
% Use as \begin{Itemize} \item ... \item ... \end{Itemize}.
\newlist{Itemize}{itemize}{1}
\setlist[Itemize]{label=--, leftmargin=1em, labelindent=0pt, labelsep=0.5em, itemsep=0.5ex, topsep=0.5ex}

\newcommand{\TODO}[2]{%
    \noindent\textcolor{gray}{%
            {\scriptsize\textbf{TODO(#1)}: #2}
    }
}

\setlength{\belowcaptionskip}{-1ex}
\setlength{\textfloatsep}{0.2ex}
\newenvironment{Figure}{
    \begin{figure}[tb]
} {
    \end{figure}
}

% Full-width variant (spans both columns in the 2-column layout).
\newenvironment{WideFigure}{
    \begin{figure*}[!t]
} {
    \end{figure*}
}

% Full-width table variant (spans both columns in the 2-column layout).
\newenvironment{WideTable}{
    \begin{table*}[!t]
} {
    \end{table*}
}

% Abbreviations

\newcommand{\MVP}[0]{\textsf{MVP}\xspace}
\newcommand{\MSL}[0]{\textsf{MSL}\xspace}
\newcommand{\WP}[0]{WP\xspace}
\newcommand{\solidity}[0]{\textsf{Solidity}\xspace}
\newcommand{\kay}[0]{\textsf{K}\xspace}
\newcommand{\fstar}[0]{\textsf{F*}\xspace}
\newcommand{\verus}[0]{\textsf{Verus}\xspace}
\newcommand{\certora}[0]{\textsf{Certora}\xspace}
\newcommand{\CC}[0]{\textsf{CC}\xspace}


% IVL

% Based on esz.sty -- see CLAUDE.md

\newenvironment{ivl}{\zedindent=2ex\begin{zed}}{\end{zed}}

\newcommand{\requiresof}[2]{\Zpreop{requires\_of}\langle#1\rangle(#2)}
\newcommand{\abortsof}[2]{\Zpreop{aborts\_of}\langle#1\rangle(#2)}
\newcommand{\ensuresof}[2]{\Zpreop{ensures\_of}\langle#1\rangle(#2)}
\newcommand{\resultof}[2]{\Zpreop{result\_of}\langle#1\rangle(#2)}
\newcommand{\modifiesof}[1]{\Zpreop{modifies\_of}\langle#1\rangle}

% Variant-set names used in the implementation section.
\newcommand{\Clo}{\mathcal{C}}
\newcommand{\Par}{\mathcal{P}}
\newcommand{\Fld}{\mathcal{F}}

\newcommand{\Addr}{\Zpreop{address}}

\newcommand{\Mem}{\mathcal{M}}
\newcommand{\Mod}{\Delta}
\newcommand{\Type}{\mathcal{T}}
\newcommand{\Expr}{\mathcal{E}}

% Keywords for  IVL pseudocode
% ... esz already has \LET \IF \THEN \ELSE
\newcommand{\PROC}{\Zkeyword{proc}}
\newcommand{\FUN}{\Zkeyword{fun}}
\newcommand{\AXIOM}{\Zkeyword{axiom}}
\newcommand{\TRIGGER}{\Zinrel{trigger}}
\newcommand{\RETURNS}{\Zinrel{returns}}
\newcommand{\HAVOC}{\Zkeyword{havoc}}
\newcommand{\ASSUME}{\Zkeyword{assume}}
\newcommand{\ASSERT}{\Zkeyword{assert}}
\newcommand{\DATATYPE}{\Zkeyword{datatype}}
\newcommand{\IS}{\Zinrel{is}}
\newcommand{\TRUE}{\Zkeyword{true}}
\newcommand{\OLD}{\Zpreop{old}}
\newcommand{\BOOL}{\Zpreop{bool}}
\newcommand{\READ}{\Zpreop{read}}
\newcommand{\CALL}{\Zpreop{call}}
\newcommand{\VAR}{\Zkeyword{var}}
\newcommand{\FOREACH}{\Zkeyword{foreach}}

% State-labeled expression: \slabeled{S}{e} renders as S \sat e.
% \sat is a bar-then-tilde composite that reads like the literal |~.
% No standard single-glyph relation in amssymb / stmaryrd matches
% (closest is \vDash, which is |=). We therefore compose: \scriptstyle
% shrinks both glyphs; -1mu tightens the tilde against the bar while
% keeping a small visible gap.
% Kept in sync with the Move listings literate mapping below.
\newcommand{\sat}{\mathrel{{\scriptstyle\vert}\mkern-1mu{\scriptstyle\sim}}}
\newcommand{\slabeled}[2]{#1 \sat #2}

\newcommand{\sq}[1]{\overline{#1}}


% Source Transformation

\newcommand{\transform}[0]{\Large{$\leadsto$}}


%%%% Code

% Keep IEEEtran's default Times font; do not override with charter.
%\usepackage[charter]{mathdesign}
%\def\rmdefault{bch}
%\def\ttdefault{blg}

\lstdefinestyle{MoveStyle}{
    basicstyle=\ttfamily,
    keywordstyle=\color{black}, % don't bold keywords which is the default
    commentstyle=\color{gray}\normalfont\itshape,
    escapechar=@, % use to embed LaTeX into code
    breaklines=true,
    % \sat is set in \scriptstyle and followed by tight kerning, so its
    % right side-bearing is small; tail the replacement with \, so the
    % symbol doesn't abut the next token visually.
    literate={|~}{{$\sat$\,}}1
        {<=}{{$\leq$}}2
        {>=}{{$\geq$}}2
        {==}{{$=$}}2
        {!=}{{$\neq$}}2
        {==>}{{$\Longrightarrow$}}3
        {forall}{{$\forall$}}1
        {exists}{{$\exists$}}1,
}


\lstdefinelanguage{Move}{
    morekeywords={
        abort,
        aborts_if,
        aborts_of,
        acquires,
        address,
        apply,
        as,
        assert,
        assume,
        borrow_global,
        borrow_global_mut,
        break,
        const,
        continue,
        copy,
        copyable,
        define,
        drop,
        else,
        ensures,
        ensures_of,
        exists,
        false,
        forall,
        friend,
        fun,
        global,
        has,
        havoc,
        if,
        in,
        include,
        invariant,
        key,
        let,
        loop,
        match,
        modifies,
        modifies_of,
        module,
        move,
        move_from,
        move_to,
        mut,
        native,
        num,
        old,
        onabort,
        pragma,
        proof,
        public,
        requires,
        requires_of,
        resource,
        result_of,
        return,
        schema,
        script,
        signer,
        spec,
        split,
        store,
        struct,
        true,
        u8,
        u64,
        u128,
        u256,
        update,
        use,
        with,
        where,
        while},
    sensitive=true,
    morecomment=[l]{//},
    morecomment=[s]{/*}{*/},
}

% Short inline delimiters for Move code: both |...| and !...! work.
% The ! alternative is useful when the snippet itself contains |,
% e.g. function-type syntax !|T1,...,Tn| -> T!.
%
% Inline snippets must never wrap mid-token (e.g. between & and T in |&T|),
% but they may still break at whitespace so a long inline snippet does not
% overflow the column. Block listings (Move environment) keep MoveStyle's
% default break-anywhere behavior.
\lstMakeShortInline[language=Move,style=MoveStyle,breaklines=true,breakatwhitespace=true,escapechar=@]|
\lstMakeShortInline[language=Move,style=MoveStyle,breaklines=true,breakatwhitespace=true,escapechar=@]!


% This defines a new environment for Move code.
\lstnewenvironment{Move}{
    \lstset{
        language=Move,
        style=MoveStyle,
        basicstyle=\relsize{-1}\ttfamily,
        keywordstyle=\color{blue},
    }
}{
}

% This defines a new environment for Move code in a box, without line numbers.
\lstnewenvironment{MoveBox}{
    \lstset{
        language=Move,
        style=MoveStyle,
        basicstyle=\relsize{-1}\ttfamily,
        keywordstyle=\color{blue},
    %numbers=left,
    %numberstyle=\scriptsize\color{gray},
    %frame=single,
    }
}{
}

% This defines a new environment for Move code in a box, with line numbers.
\lstnewenvironment{MoveBoxNumbered}{
    \lstset{
        language=Move,
        style=MoveStyle,
        basicstyle=\relsize{-1}\ttfamily,
        keywordstyle=\color{blue},
        numbers=left,
        numberstyle=\scriptsize\color{gray},
    %frame=single,
    }
}{
}

% This defines a new environment for diagnostics as produced by the prover.
\lstnewenvironment{MoveDiag}{
    \lstset{
        style=MoveStyle,
        basicstyle=\relsize{-2}\ttfamily,
    }
}{
}

% Environment for Claude Code transcripts.
%
% Source convention: lines that begin with "> " are the user prompt;
% everything else is the model's response. The "> " is rendered as "❯ "
% in the PDF (via a literate substitution) so the output matches the
% Claude Code CLI's prompt glyph without forcing the author to type a
% non-ASCII character.
%
% Prompt lines are styled bold and blue so they stand out; response lines
% are rendered in plain monospace, mimicking how the transcript looks in
% a terminal.
%
% Unicode glyphs commonly emitted by Claude Code (tree connectors,
% checkboxes, ellipsis, arrows, check/cross marks) are mapped via the
% literate= option so they render correctly under pdfLaTeX (where listings
% does not handle multi-byte input directly).
%
% Wrapping: long prompts wrap with breakindent matched to the width of
% "❯ " in monospace, so continuation lines line up under the prompt text
% (approximately --- listings uses a fixed indent, not a measured prefix).
%
% Limitation: ">" at the start of a line is reserved as the prompt marker.
% A literal ">" inside a response line must be escaped via the listings
% escapechar (@), e.g. @\textgreater @.
\lstnewenvironment{claude}{
    \lstset{
        basicstyle=\relsize{-1}\ttfamily,
        breaklines=true,
        breakautoindent=true,
        breakatwhitespace=true,
        breakindent=1.2em,
        columns=fullflexible,
        keepspaces=true,
        escapechar=@,
        extendedchars=true,
        inputencoding=utf8,
        literate=
            {⎿}{{$\llcorner$}}1
            {└}{{$\llcorner$}}1
            {├}{{$\vdash$}}1
            {─}{{-}}1
            {│}{{$\vert$}}1
            {◻}{{$\square$}}1
            {◼}{{$\blacksquare$}}1
            {☐}{{$\square$}}1
            {☑}{{$\boxtimes$}}1
            {●}{{$\bullet$}}1
            {○}{{$\circ$}}1
            {◯}{{$\bigcirc$}}1
            {⏺}{{$\bullet$}}1
            {✽}{{$\ast$}}1
            {★}{{$\star$}}1
            {☆}{{$\star$}}1
            {…}{{\ldots}}1
            {↗}{{$\nearrow$}}1
            {↘}{{$\searrow$}}1
            {↳}{{$\hookrightarrow$}}1
            {✓}{{\Checkmark}}1
            {✗}{{\XSolidBrush}}1
            {•}{{$\bullet$}}1
            {–}{{--}}1
            {—}{{---}}1
            {‘}{{`}}1
            {’}{{'}}1
            {“}{{``}}1
            {”}{{''}}1,
        moredelim=[l][\color{blue}\bfseries]{>},
    }
}{
}


%%% Local Variables:
%%% mode: latex
%%% TeX-master: "main"
%%% End:
